\documentclass[12pt]{article}
\usepackage{amsmath, amssymb, amsthm}
\usepackage{geometry}
\usepackage{hyperref}
 
\geometry{margin=1in}
 
\newtheorem*{theorem}{Theorem}
\theoremstyle{remark}
\newtheorem*{note}{Key Concept}
 
\newcommand{\Lc}{\mathcal{L}}
 
\title{\textbf{Axler Exercise 3D.1: The Inverse of an Inverse}\\
\large A Proof and Conceptual Commentary}
\author{Linear Algebra Done Right, 4th Edition}
\date{}
 
\begin{document}
 
\maketitle
 
\section*{Context in the Text}
 
This is Exercise 1 in Section 3D (\emph{Invertibility and Isomorphisms}) of \emph{Linear Algebra Done Right}, 4th edition. Section 3D opens by
defining invertibility and the inverse of a linear map (Definition 3.59),
establishing that inverses are unique (Theorem 3.60), and introducing the
notation $T^{-1}$ (3.61). This exercise is the first in the section and is
intentionally elementary: it asks the reader to check that the defining
equations of an inverse are symmetric in $T$ and $T^{-1}$, and to invoke
uniqueness. It is less a proof than a careful reading of the definition.
 
\section*{Prerequisites}
 
\begin{itemize}
    \item \textbf{Definition 3.59 (Invertible, inverse):} $T \in \Lc(V, W)$ is
          invertible if there exists $S \in \Lc(W, V)$ such that $ST = I_V$ and
          $TS = I_W$. Such an $S$ is called an inverse of $T$.
 
    \item \textbf{Theorem 3.60 (Uniqueness of the inverse):} An invertible
          linear map has a unique inverse, denoted $T^{-1}$.
 
    \item \textbf{Notation 3.61:} If $T \in \Lc(V,W)$ is invertible, then
          $T^{-1}$ is the unique element of $\Lc(W, V)$ satisfying
          $T^{-1}T = I_V$ and $TT^{-1} = I_W$.
\end{itemize}
 
\section*{Statement}
 
\begin{theorem}[Exercise 3D.1]
Suppose $T \in \Lc(V, W)$ is invertible. Then $T^{-1}$ is invertible and
\[
    \left(T^{-1}\right)^{-1} = T.
\]
\end{theorem}
 
\section*{Proof}
 
Since $T$ is invertible, its inverse $T^{-1} \in \Lc(W, V)$ satisfies
\[
    T^{-1}T = I_V \qquad \text{and} \qquad TT^{-1} = I_W.
\]
These same two equations, read from the perspective of $T^{-1}$, say precisely
that $T \in \Lc(V, W)$ is an inverse of $T^{-1} \in \Lc(W, V)$: a map in
$\Lc(V, W)$ that composes with $T^{-1}$ on each side to give the appropriate
identity. Therefore $T^{-1}$ is invertible.
 
Since inverses are unique (Theorem 3.60), $T$ must be \emph{the} inverse of
$T^{-1}$, i.e.,
\[
    \left(T^{-1}\right)^{-1} = T. \qquad \blacksquare
\]
 
\section*{Key Concept}
 
\begin{note}
The core observation is that Definition 3.59 is \emph{symmetric} in $T$ and
$T^{-1}$. The two conditions
\[
    T^{-1}T = I_V \qquad \text{and} \qquad TT^{-1} = I_W
\]
make no distinction between which map is the ``original'' and which is the
``inverse'' --- each is the inverse of the other. The only thing the proof needs
to add beyond this symmetry observation is uniqueness (Theorem 3.60), to
upgrade ``$T$ is \emph{an} inverse of $T^{-1}$'' to ``$T$ is \emph{the}
inverse of $T^{-1}$,'' which is what the notation $(T^{-1})^{-1} = T$ asserts.
 
This same symmetry argument recurs throughout algebra: in group theory, the
inverse of $g^{-1}$ is $g$ for exactly the same reason.
\end{note}
 
\end{document}
